% D5 CJ-first boundary reset proof reduction 056
\begin{proposition}[Carry\_jump boundary reduction]
Fix an odd modulus \(m\) and the active best-seed branch for \(R1\to H_{L1}\).
Assume the theorem-side coordinates
\[
B=(s,u,v,\lambda,f),\qquad \tau,\qquad \epsilon_4.
\]
On the boundary \(\tau=0\) with \(\epsilon_4=\mathrm{carry\_jump}\), the checked reset law
through \(m\le 23\) is
\[
R_{\mathrm{cj}}(s,v,\lambda)=
\begin{cases}
0,& s+v+\lambda\equiv 2 \pmod m,\\
1,& s=1 \text{ and } s+v+\lambda\not\equiv 2 \pmod m,\\
m-2,& \text{otherwise.}
\end{cases}
\]
Equivalently, the current raw coordinate satisfies
\[
q\equiv 1-s-v-\lambda \pmod m.
\]
Moreover, if \(\rho=u_{\mathrm{source}}+1\) is the constructive source residue from the
compute-side refinement, then on the carry\_jump boundary
\[
q\equiv u-\rho+1 \pmod m,
\qquad
\rho\equiv s+u+v+\lambda \pmod m.
\]
Therefore the uniform carry\_jump proof reduces to the noncarry dichotomy
\[
q\neq m-1
\Longrightarrow
R_{\mathrm{cj}}=
\begin{cases}
1,& s=1,\\
m-2,& s\neq 1.
\end{cases}
\]
\end{proposition}

\begin{remark}[Proof program]
The boundary reset theorem now splits into three parts:
\begin{enumerate}
\item wrap \(\mapsto 0\), already formal;
\item the carry\_jump lemma above, whose only remaining content is the noncarry dichotomy;
\item the \texttt{other} branch lemma, which still requires the two-subtype reduction.
\end{enumerate}
Thus the proof burden is no longer ``decode the future window'' but ``prove two explicit
boundary branch lemmas, starting with carry\_jump.''
\end{remark}
